\section{SMART goal 3: Investigation of data lineage for metrics}
\subsection*{Goal definition}
\textbf{3) Investigation of data lineage for metrics}

\textbf{WHY:} To apply DQM checks to the correct data sources and enable root cause analysis.

By the end of the analysis phase, investigate and document the data lineage from data ingestion to the reporting layer.

Success is measured by a documented lineage description that allows data to be followed through the system and its transformations to be understood, including the validation possibilities available to responsible stakeholders.



\subsection*{Core results}
\begin{itemize}
    
\item \textbf{Business report documentation:} The business report for message processing DQM was documented. The key information was identified as \textbf{Diff\_Discharged\_Submitted}. Both submitted and discharged values should be in the same number range, as this check serves to ensure data completeness. This is because every mail piece in a sorting center has one submitted and one discharged event count per day, as the machines are empty at the end of every day. Therefore, the difference between these values needs to be below the threshold of 0.01. This is described in section \ref{Implementation1SubmittedDischargedDifference}.

\item \textbf{Key value identification:} In the implementation chapter, the quality metrics of submitted and discharged, in conjunction with the difference between the two (\textbf{Diff\_Discharged\_Submitted}), were identified. These metrics represent the key values used to evaluate the daily completeness of data (Section \ref{Implementation1SubmittedDischargedDifference}).

\item \textbf{SAC report documentation:} The report information relevant to the MVP was identified, including the date, three-letter location key, submitted number, discharged number, and their ratio, and was documented in section \ref{sectionSVADQMReport}.

\item \textbf{DSP data model description:} Within the DSP, the data lineage of the entire report was examined. A key finding was the location where the ratio is calculated, namely within the DSP data model. Furthermore, it was identified that the data for the MVP case is based solely on AWS Athena; see section \ref{ImplementationInvestigateDSP}.

\item \textbf{Business stakeholder confirmation:} The AWS Athena query was traced, and during a meeting with the stakeholders, the key information was confirmed. This is described in section \ref{ImplementationFourBusinessMeeting}.

\item \textbf{AWS Athena query:} A concise analysis of the query structure and the AWS Athena–specific SQL was conducted, identifying the first relevant components. Finally, documentation was created and presented to the DQM team, as described in section \ref{subsectionAthenaQuery}.

\item \textbf{DQM team update:} To update the DQM team, the data lineage from the AWS Athena query, including relevant accounts, was documented. This includes a brief introduction to the use case and the query for the MVP in section \ref{ImplementationFourConfluenceDocu}. The meeting is documented in appendix \ref{apendixDQM1911}.

\end{itemize}

\subsection*{Supporting results}

\begin{itemize}
\item \textbf{Data lineage description:} The data lineage was described from: \textit{1) the physical data source:} scans are performed on parcels by different scanners, \textit{2) the initial data storage:} the scan events are first saved in AWS S3, \textit{3) Kafka ingestion:} Apache Kafka ingests the events into a topic, \textit{4) data delivery to DWH:} Apache Kafka is consumed by DWH LS OPS, where the data is saved to the raw layer. This is described in Section \ref{SolutionDataLineage}.

\item \textbf{Reporting pipeline description:} The reporting pipeline of DWH LS OPS is defined as the \textbf{storage layer} using Amazon S3, followed by the \textbf{query and processing layer} using Amazon Athena and SAP Datasphere (DSP), and finally the \textbf{consumption layer} using SAP Analytics Cloud (SAC). These steps are detailed in Section \ref{SolutionReportingPipeline}.

\item \textbf{Placement of data quality:} Data quality checks are best placed on the processing and transformation layer. This allows both meaningful and effective tests. Earlier checks on the data lineage and reporting pipeline are less effective, as only global data completeness can be verified, while later checks on the reporting layer remove the possibility of corrective action. This is described in Section \ref{ResearchDataQChecksOnDataLifecycle}.

\item \textbf{Key value identification:} In the implementation chapter, the quality metrics of submitted and discharged events, in conjunction with the difference between the two (\textbf{Diff\_Discharged\_Submitted}), are identified. These metrics represent the key values used to evaluate the daily completeness of data (Section \ref{Implementation1SubmittedDischargedDifference}).
\end{itemize}

\subsection*{What was not achieved?}

The goal was completely fulfilled; nothing was left unachieved.


\subsection*{Goal verdict}
This goal was completely fulfilled. All required deliverables have been provided for this SMART goal.
These include:
\begin{multicols}{2}
\begin{itemize}
    \item \textbf{Supporting} Data lineage description
    \item \textbf{Supporting} Reporting pipeline description
    \item \textbf{Supporting} Placement of data quality
    \item \textbf{Supporting} Key value identification    
    \item \textbf{Core} Business report documentation
    \item \textbf{Core} Key value identification
    \item \textbf{Core} SAC report documentation
    \item \textbf{Core} DSP data model description
    \item \textbf{Core} Business stakeholder confirmation
    \item \textbf{Core} AWS Athena query 
    \item \textbf{Core} DQM team update
\end{itemize}
\end{multicols}


\subsection*{Learning and outlook for future work}

This allowed the identification of the key measurements for DQM. As this SMART goal was fully successful, no further learnings are carried forward to subsequent projects, aside from the tools SAC and DSP, which were examined in detail during their fulfillment.